%============================================================================
% LICENSING NOTICE
%============================================================================
% This WIP file, upon integration into guide-v.tex, becomes part of the
% TMS/DMS/CMS Usage Guide and is released under CC BY-NC 4.0.
% For details, see LICENSE in the repository root.
%============================================================================

%============================================================================
% FALCON BMS TMS/DMS/CMS HOTAS GUIDE
% WIP FILE TEMPLATE — FINAL (Briefing)
%============================================================================
% IMPORTANTE: Este é um TEMPLATE padronizado para arquivos WIP (Work-In-Progress)
% que serão integrados ao guide.tex.
% Nomenclatura: Siga wip-naming
% Padrão: section-C{N}-S{M}[-S{K}]-{titulo}-{status}-{data}.tex
% Exemplo: section-C5-S2-cms-actuation-hotas-tables-final-2026-01-10.tex
% Status: dev | review | final | alpha | approved | deprecated
% Note: 'alpha' is the parallel pre-release flow (optional) — see wip-naming.md §0.5
% Locação: WIP/ (ativo) | ARCHIVE/ (aprovado/descartado)

%============================================================================
% PREAMBLE COMPLETO — TMS/DMS/CMS Usage Guide for Falcon BMS 4.38.1
% Gerado: 25 March 2026
%============================================================================

\documentclass[11pt, a4paper, twoside]{report}

% --------------------------------------------------------------------------
% BASIC ENCODING AND LANGUAGE
% --------------------------------------------------------------------------

\usepackage[utf8]{inputenc}
\usepackage[T1]{fontenc}
\usepackage[english]{babel}

% --------------------------------------------------------------------------
% FONTS AND MICROTYPOGRAPHY
% --------------------------------------------------------------------------

\usepackage{lmodern}
\usepackage{microtype}

% --------------------------------------------------------------------------
% PAGE GEOMETRY AND LAYOUT
% --------------------------------------------------------------------------

\usepackage{geometry}
\geometry{a4paper, left=2.0cm, right=2.0cm, top=2.5cm, bottom=2.5cm}
\usepackage{setspace}
\onehalfspacing

% --------------------------------------------------------------------------
% VERSION CONTROL MACROS
% --------------------------------------------------------------------------

\newcommand{\docversion}{v0.4.2.0-alpha.3}
\newcommand{\docbuild}{2026-03-30}
\newcommand{\docstartdate}{05 January 2026}
\newcommand{\docenddate}{20260327}
\newcommand{\chapterscompletedof}{3/6}
\newcommand{\tablesfilledpct}{Chapter 5}
\newcommand{\fulldocversion}{\docversion+\docbuild}

% --------------------------------------------------------------------------
% COLORS AND LINKS AND PDF SETTINGS
% --------------------------------------------------------------------------

\usepackage[table]{xcolor}
\definecolor{linkblue}{HTML}{004488}
\definecolor{linkred}{HTML}{882222}
\definecolor{headerblue}{HTML}{003366}
\definecolor{rowgray}{HTML}{F5F5F5}
\definecolor{subheadgray}{HTML}{E0E0E0}
\usepackage{soul}
\usepackage[pdfencoding=auto, psdextra, colorlinks=true, linkcolor=linkblue, citecolor=linkred, urlcolor=linkblue, breaklinks=true]{hyperref}
\usepackage{hyperxmp}
\hypersetup{
	pdftitle={TMS, DMS and CMS Usage Guide for Falcon BMS},
	pdfauthor={Carlos "Metal" Nader},
	pdfsubject={Flight Simulation - Falcon BMS HOTAS Reference},
	pdfkeywords={Falcon BMS, F-16, HOTAS, TMS, DMS, CMS, Flight Simulation},
	pdfcreator={pdfLaTeX (TeXstudio)},
	pdfproducer={TeXstudio},
	pdflang={en-US},
	pdfcopyright={CC BY-NC 4.0 · Carlos "Metal" Nader},
	pdflicenseurl={https://creativecommons.org/licenses/by-nc/4.0/},
	pdfcontacturl={https://github.com/carlos-nader/tms-dms-cms-usage-guide},
	pdfdisplaydoctitle=true,
	bookmarksnumbered=true,
	bookmarksopen=true,
	bookmarksdepth=1,
	pdfstartview=FitH,
}
\pdfinfo{/CreationDate (D:\docbuild) /ModDate (D:\docbuild)}
\usepackage{bookmark}

% --------------------------------------------------------------------------
% CAPTION SETUP BY TYPE
% --------------------------------------------------------------------------

\usepackage{caption}

% GLOBAL PATTERN
\captionsetup{font=small, labelfont=bf, justification=centering, singlelinecheck=true}
% FIGURES
\captionsetup[figure]{font=footnotesize, labelfont=bf, justification=centering, singlelinecheck=true}
% % TABLE/LONGTABLE/hotastable:
\captionsetup[table]{font=small, labelfont=bf, justification=centering, singlelinecheck=true}

% --------------------------------------------------------------------------
% HEADERS AND FOOTERS (IMPROVED for report + twoside)
% --------------------------------------------------------------------------

\usepackage{fancyhdr}
\setlength{\headheight}{25pt}
\pagestyle{fancy}
\fancyhf{}
\fancyhead[LO,RE]{\small\textit{\leftmark}}
\fancyhead[RO,LE]{\small\thepage}
\fancyfoot{}
\renewcommand{\headrulewidth}{0.4pt}
\renewcommand{\footrulewidth}{0pt}

% --------------------------------------------------------------------------
% CHAPTER FORMATTING AND SPACING (via titlesec)
% --------------------------------------------------------------------------

\usepackage{titlesec}

% --------------------------------------------------------------------------
% CHAPTER - Nível 0 (Destaque Máximo)
% --------------------------------------------------------------------------
% Shape: display → standalone em nova linha
% Font: \Large\bfseries --- label e título ambos em bold para destaque máximo
% Numeração: SIM
% Espaçamento: 15pt antes, 28pt depois
% --------------------------------------------------------------------------

\titleformat{\chapter}[display]
{\normalfont\Large\bfseries}           % Formato: LARGE + bold (label)
{\chaptertitlename~\thechapter}        % Rótulo: "Chapter 1" (bold também)
{20pt}                                 % Espaço vertical antes do corpo
{\Large\bfseries}                      % Corpo do título: LARGE + bold

\titlespacing{\chapter}
{0pt}       % Margem esquerda (sem indent)
{15pt}      % Espaço ANTES do título do capítulo
{28pt}      % Espaço DEPOIS do título do capítulo (AUMENTADO)
[0pt]       % Margem direita

% --------------------------------------------------------------------------
% SECTION - Nível 1
% --------------------------------------------------------------------------
% Shape: hang → label à esquerda, corpo indentado (compacto)
% Font: \large\bfseries (mesmo destaque que chapter para coerência)
% Numeração: SIM
% Espaçamento: 15pt antes, 8pt depois (REDUZIDO de 15pt/10pt anteriores)
% --------------------------------------------------------------------------

\titleformat{\section}[hang]
{\normalfont\large\bfseries}           % Formato: large + bold
{\thesection}                          % Rótulo: "1", "2", etc
{1em}                                  % Espaço entre rótulo e corpo
{}                                     % Corpo do título (herda do format)

\titlespacing{\section}
{0pt}       % Margem esquerda (sem indent)
{30pt}      % Espaço ANTES do título da seção
{8pt}       % Espaço DEPOIS do título da seção
[0pt]       % Margem direita

% --------------------------------------------------------------------------
% SUBSECTION - Nível 2
% --------------------------------------------------------------------------
% Shape: hang → label esquerda, corpo indentado
% Font: \normalsize\bfseries (redução visual, ainda robusto)
% Numeração: SIM
% Espaçamento: 12pt antes, 6pt depois (compacto)
% --------------------------------------------------------------------------

\titleformat{\subsection}[hang]
{\normalfont\normalsize\bfseries}      % Formato: tamanho normal + bold
{\thesubsection}                       % Rótulo: "1.1", "1.2", etc
{1em}                                  % Espaço entre rótulo e corpo
{}                                     % Corpo do título

\titlespacing{\subsection}
{0pt}       % Margem esquerda (sem indent)
{24pt}      % Espaço ANTES do título da subseção
{6pt}       % Espaço DEPOIS do título da subseção
[0pt]       % Margem direita

% --------------------------------------------------------------------------
% SUBSUBSECTION - Nível 3
% --------------------------------------------------------------------------
% Shape: hang → label esquerda, máximo recuo (bem compacto)
% Font: \normalsize\bfseries (mesmo tamanho que subsection, padrão)
% Numeração: SIM (configurable com secnumdepth)
% Espaçamento: 10pt antes, 4pt depois (bem comprimido)
% --------------------------------------------------------------------------

\titleformat{\subsubsection}[hang]
{\normalfont\normalsize\bfseries}      % Formato: tamanho normal + bold
{\thesubsubsection}                    % Rótulo: "1.1.1", "1.1.2", etc
{1em}                                  % Espaço entre rótulo e corpo
{}                                     % Corpo do título

\titlespacing{\subsubsection}
{0pt}       % Margem esquerda (sem indent)
{18pt}       % Espaço ANTES do título
{4pt}       % Espaço DEPOIS do título (bem comprimido)
[0pt]       % Margem direita

% --------------------------------------------------------------------------
% PARAGRAPH - Nível 4 (CRÍTICO: Ilusão Óptica Resolvida)
% --------------------------------------------------------------------------
% Shape: runin → inline, integrado ao texto (seu requisito)
% Font: \small\bfseries (SOLUÇÃO PARA ILUSÃO ÓPTICA DO NEGRITO)%
% Numeração: NÃO (padrão para \paragraph)
% Espaçamento: runin (sem espaço vertical antes, integrado ao texto)
% --------------------------------------------------------------------------

\titleformat{\paragraph}[runin]
{\normalfont\small\bfseries}           % Formato: SMALL + bold (compensa ilusão)
{}                                     % Rótulo: vazio (sem numeração)
{0em}                                  % Espaço entre rótulo e corpo (zero, runin)
{}                                     % Corpo do título

\titlespacing{\paragraph}
{0pt}       % Margem esquerda (sem indent)
{8pt}       % Espaço ANTES do título do parágrafo (ADICIONADO - cria separação)
{1em}       % Espaço DEPOIS do título (ADICIONADO - espaço após ":")
[0pt]       % Margem direita

% --------------------------------------------------------------------------
% SUBPARAGRAPH - Nível 5 (Preparação Futura)
% --------------------------------------------------------------------------
% Shape: runin → inline, integrado ao texto
% Font: \small (MENOS destaque que \paragraph, sem bold)
% Numeração: NÃO (padrão para \subparagraph)
% Espaçamento: runin (integrado)
% Uso: Se usado, aparecerá com destaque menor que \paragraph
% --------------------------------------------------------------------------

\titleformat{\subparagraph}[runin]
{\normalfont\small}                    % Formato: small, SEM bold (menos destaque)
{}                                     % Rótulo: vazio (sem numeração)
{0em}                                  % Espaço entre rótulo e corpo (zero, runin)
{}                                     % Corpo do título

\titlespacing{\subparagraph}
{0pt}       % Margem esquerda (sem indent)
{8pt}       % Espaço ANTES do título (ADICIONADO)
{1em}       % Espaço DEPOIS do título (ADICIONADO)
[0pt]       % Margem direita

% --------------------------------------------------------------------------
% TABLES AND MACROS
% --------------------------------------------------------------------------

\usepackage{booktabs}
\usepackage{array}
\usepackage{longtable}
\usepackage{tabularx}

% Custom Columns
\newcolumntype{L}[1]{>{\raggedright\arraybackslash}p{#1}}
\newcolumntype{C}[1]{>{\centering\arraybackslash}p{#1}}
\newcolumntype{R}[1]{>{\raggedleft\arraybackslash}p{#1}}

% Macro for Visual Reference Links
\newcommand{\imglink}[1]{\hspace{2pt}\hyperref[#1]{\scriptsize\textbf{[Fig]}}}

% --------------------------------------------------------------------------
% HOTAS LOT FILE HANDLE
% --------------------------------------------------------------------------
\makeatletter
\newcommand{\listofhotastables}{%
	\section*{List of HOTAS Tables}%
	\addcontentsline{toc}{section}{List of HOTAS Tables}%
	\@starttoc{hotas}%
}
\makeatother

% --------------------------------------------------------------------------
% HOTAS table environment
% Version: 2.1 (2026-02-10)
% --------------------------------------------------------------------------

\newenvironment{hotastable}[1]{%
	\footnotesize
	\setlength{\tabcolsep}{3pt}
	\renewcommand{\arraystretch}{1.35}
	\begin{longtable}{L{1.00cm} L{0.90cm} L{0.90cm} L{3.30cm} L{6.40cm} L{1.40cm} L{1.60cm}}
		\caption{#1}\label{table.\thetable}\\
		\noalign{\addtocontents{hotas}{\protect\contentsline{table}{\protect\numberline{\thetable}#1}{\thepage}{table.\thetable}}}%
		\rowcolor{headerblue}
		\multicolumn{1}{>{\centering\arraybackslash}p{1.00cm}}{\textbf{\color{white}Mode}} &
		\multicolumn{1}{>{\centering\arraybackslash}p{0.90cm}}{\textbf{\color{white}Dir.}} &
		\multicolumn{1}{>{\centering\arraybackslash}p{0.90cm}}{\textbf{\color{white}Act.}} &
		\multicolumn{1}{>{\centering\arraybackslash}p{3.30cm}}{\textbf{\color{white}Function}} &
		\multicolumn{1}{>{\centering\arraybackslash}p{6.40cm}}{\textbf{\color{white}Effect / Nuance}} &
		\multicolumn{1}{>{\centering\arraybackslash}p{1.40cm}}{\textbf{\color{white}Dash34}} &
		\multicolumn{1}{>{\centering\arraybackslash}p{1.60cm}}{\textbf{\color{white}Train.}} \\
		\endfirsthead
		\rowcolor{headerblue}
		\multicolumn{1}{>{\centering\arraybackslash}p{1.00cm}}{\textbf{\color{white}Mode}} &
		\multicolumn{1}{>{\centering\arraybackslash}p{0.90cm}}{\textbf{\color{white}Dir.}} &
		\multicolumn{1}{>{\centering\arraybackslash}p{0.90cm}}{\textbf{\color{white}Act.}} &
		\multicolumn{1}{>{\centering\arraybackslash}p{3.30cm}}{\textbf{\color{white}Function}} &
		\multicolumn{1}{>{\centering\arraybackslash}p{6.40cm}}{\textbf{\color{white}Effect / Nuance}} &
		\multicolumn{1}{>{\centering\arraybackslash}p{1.40cm}}{\textbf{\color{white}Dash34}} &
		\multicolumn{1}{>{\centering\arraybackslash}p{1.60cm}}{\textbf{\color{white}Train.}} \\
		\endhead
		\multicolumn{7}{r}{\small\emph{Continued on next page}}\\
		\endfoot
		\endlastfoot
	}{%
	\end{longtable}
}

% --------------------------------------------------------------------------
% SIMPLE REFERENCE MACROS FOR BMS DOCS
% --------------------------------------------------------------------------

\newcommand{\dashref}[1]{Dash-34~\S~#1}
\newcommand{\dashone}[1]{Dash-1~\S~#1}
\newcommand{\trnref}[1]{TRN~#1}
\newcommand{\trnman}{BMS Training Manual 4.38.1}
\newcommand{\bmsver}{Falcon BMS~4.38.1}
\newcommand{\dashrefs}[1]{\textit{TO 1F-16CMAM-34-1-1}, Dash-34, sections \texttt{#1}}

% --------------------------------------------------------------------------
% CROSS-REFERENCE MACROS (INTERNAL GUIDE)
% --------------------------------------------------------------------------

\newcommand{\secref}[1]{\hyperref[#1]{Section~\ref*{#1}}}
\newcommand{\chapref}[1]{\hyperref[#1]{Chapter~\ref*{#1}}}
\newcommand{\tabref}[1]{\hyperref[#1]{Table~\ref*{#1}}}
\newcommand{\figref}[1]{\hyperref[#1]{Figure~\ref*{#1}}}

% --------------------------------------------------------------------------
% GRAPHICS
% --------------------------------------------------------------------------

\usepackage{graphicx}
\graphicspath{{fig/}}
\usepackage{float}
\usepackage{enumitem}

% --------------------------------------------------------------------------
% TITLE
% --------------------------------------------------------------------------

\title{TMS, DMS and CMS Usage Guide for \bmsver}
\author{Carlos ``Metal'' Nader}
\date{Version \fulldocversion{} | Progress: Chapters \chapterscompletedof{} | Tables \tablesfilledpct{} | January 2026}

% ============================================================================
% DOCUMENT BEGIN
% ============================================================================

\begin{document}

	\maketitle

	\pagenumbering{roman}

	% --------------------------------------------------------------------------
	% TOC DEPTH CONFIGURATION
	% --------------------------------------------------------------------------
	\setcounter{tocdepth}{3}       % Show up to \subsubsection in TOC
	\setcounter{secnumdepth}{3}    % Number up to \subsubsection

	\newpage
	\tableofcontents
	\newpage

	% --------------------------------------------------------------------------
	% LIST OF HOTAS TABLES
	% --------------------------------------------------------------------------
	\phantomsection
	\listofhotastables
	\newpage

	\pagenumbering{arabic}

%============================================================================
% CONTENT BEGINS HERE (use \chapter{}, \section{}, \subsection{}, etc.)
%============================================================================

%--------------------------------------------------------------------------
% CHAPTER 1: INTRODUCTION
%--------------------------------------------------------------------------
\chapter{Introduction}
\label{chap:C1}

This document is a community-made reference guide for \bmsver, focused on practical use of three HOTAS controls: the Target Management Switch (TMS), the Display Management Switch (DMS), and the Countermeasures Management Switch (CMS). The fundamental behavior of these switches has remained consistent since at least Falcon BMS 4.36, making this guide applicable to virtually any Falcon BMS user.

This work is entirely unofficial. The author is not affiliated with Benchmark Sims, MicroProse, any real-world air force, or any aircraft or weapons manufacturer. All interpretations, simplifications, errors, and omissions are solely the responsibility of the author and must not be attributed to the Falcon BMS development team or any real-world organization.

\textbf{Nothing in this document should ever be used for real-world operations, training, or procedures.}

%--------------------------------------------------------------------------
% SECTION 1.1: SOURCES AND REFERENCES
%--------------------------------------------------------------------------
\section{Sources and references}
\label{sec:C1-S1}

This guide is grounded in the following primary Falcon BMS documents:

\begin{enumerate}
	\item \textbf{TO BMS 1F-16CMAM-34-1-1} (Dash-34, Change 4.38) --- Avionics and Nonnuclear Weapons Delivery Flight Manual. Primary source for HOTAS behavior, SOI logic, and weapons employment.
	\item \textbf{BMS Training Manual 4.38.1} (October 2025) --- Training missions and learning objectives. Used for practical cross-references and recommended training progression.
\end{enumerate}

All technical claims in this guide are cross-referenced to these sources. Where specific sections are cited, the notation \dashref{X.Y.Z} refers to Dash-34 section X.Y.Z.

%--------------------------------------------------------------------------
% SECTION 1.2: SCOPE AND PURPOSE
%--------------------------------------------------------------------------
\section{Scope and purpose}
\label{sec:C1-S2}

The guide organizes TMS, DMS, and CMS behavior by operational context: master mode, sensor configuration, and weapon employment scenario, presenting each switch input and its effect in mode-based tables. Other systems, avionics, and controls are described only when essential to understanding the behavior of these three switches.

Knowledge of basic F-16 operation and familiarity with master modes (NAV, A-A, and A-G) is assumed. This is not a comprehensive HOTAS or avionics manual and does not replace the Dash-34 or Training Manual.

%--------------------------------------------------------------------------
% SECTION 1.3: DOCUMENT STRUCTURE AND HOW TO READ IT
%--------------------------------------------------------------------------
\section{Document structure and how to read it}
\label{sec:C1-S3}

%--------------------------------------------------------------------------
% SECTION 1.3.1: CHAPTER OVERVIEW
%--------------------------------------------------------------------------
\subsection{Chapter overview}
\label{sec:C1-S3-S1}

This guide is organized into six chapters:

\begin{itemize}
	\item \textbf{Chapter 2} establishes the conceptual framework for all three switches: Sensor of Interest (SOI), master modes, and a high-level overview of DMS, TMS, and CMS.
	\item \textbf{Chapters 3--5} cover one switch each --- DMS (Chapter 3), TMS (Chapter 4), and CMS (Chapter 5) --- combining narrative explanation with standardized hotastable tables organized by master mode and sensor configuration.
	\item \textbf{Chapter 6} is a quick reference chapter, combining the hotastable tables from Chapters 3--5 with switch diagrams for immediate visual lookup.
\end{itemize}

%--------------------------------------------------------------------------
% SECTION 1.3.2: TABLE FORMAT (HOTASTABLE)
%--------------------------------------------------------------------------
\subsection{Table format}
\label{sec:C1-S3-S2}

Chapters 3, 4, and 5 use a standardized seven-column table format called \texttt{hotastable} to document switch behavior. Each row describes a single switch action in a specific operational context:

\begin{table}[H]
	\centering
	\caption{HOTAS Table Column Definitions}
	\label{tab:C1-hotastable-format}
	\small
	\renewcommand{\arraystretch}{1.25}
	\begin{tabular}{L{2.2cm} L{11.5cm}}
		\toprule
		\textbf{Column} & \textbf{Description} \\
		\midrule
		Mode & Operational context: master mode, sensor state, or weapon configuration \\
		Dir. & Switch direction: Up, Down, Left, or Right \\
		Act. & Actuation type: Short press or Long press \\
		Function & Brief name of the function performed \\
		Effect / Nuance & Detailed explanation of system behavior, constraints, and tactical considerations \\
		Dash34 & Cross-reference to Dash-34 section(s) \\
		Train. & Recommended BMS training mission number(s) \\
		\bottomrule
	\end{tabular}
\end{table}

%--------------------------------------------------------------------------
% SECTION 1.3.3: TYPOGRAPHIC CONVENTIONS
%--------------------------------------------------------------------------
\subsection{Typographic conventions}
\label{sec:C1-S3-S3}

Three inline formatting conventions are used consistently throughout the guide:

\begin{itemize}
	\item \textbf{Bold} identifies the three HOTAS switches covered by this guide and their directions when referenced in prose: \textbf{TMS Up}, \textbf{CMS Down}, \textbf{DMS Left}.
	\item \texttt{Monospace} reproduces strings exactly as they appear on a system display or are generated as system voice messages: \texttt{DISPENSE RDY}, \texttt{COUNTER}, \texttt{STBY}, \texttt{AVNC}.
	\item \textit{Italic} marks a technical distinction where sentence structure alone does not sufficiently convey the critical information.
\end{itemize}

%--------------------------------------------------------------------------
% SECTION 1.4: ABOUT THIS DOCUMENT
%--------------------------------------------------------------------------
\section{About this document}
\label{sec:C1-S4}

%--------------------------------------------------------------------------
% SUBSECTION 1.4.1: DOCUMENT VERSION AND STATUS
%--------------------------------------------------------------------------
\subsection{Document version and status}
\label{sec:C1-S4-S1}

\begin{table}[H]
	\centering
	\caption{Document Status}
	\label{tab:C1-doc-status}
	\small
	\renewcommand{\arraystretch}{1.25}
	\begin{tabular}{L{4.5cm} L{9.0cm}}
		\toprule
		\textbf{Item} & \textbf{Value} \\
		\midrule
		Document Version & \fulldocversion \\
		Falcon BMS Version & 4.38.1 (Update 1) \\
		Development Phase & Pre-publication (0.x.x.x) \\
		Chapters with Content & \chapterscompletedof \\
		Development Start & \docstartdate \\
		Last Update & \docenddate \\
		\bottomrule
	\end{tabular}
\end{table}

%--------------------------------------------------------------------------
% SUBSECTION 1.4.2: AUTHORSHIP AND AI ASSISTANCE
%--------------------------------------------------------------------------
\subsection{Authorship and AI assistance}
\label{sec:C1-S4-S2}

This guide was created by a member of the Falcon BMS community with structured assistance from AI language models:

\begin{itemize}
	\item \textbf{Perplexity Pro} (Perplexity AI) --- Research organization, project structuring, governance framework design, coding, initial draft content generation, and technical review; used via web interface.
	\item \textbf{Claude Pro} (Anthropic) --- Research, coding, initial draft content generation, technical review, and governance framework; used via web interface and VSCode (Claude Code).
	\item \textbf{GitHub Copilot Pro} (GitHub / Microsoft) --- Documentation, coding, repository management, issue and pull request coordination, and collaborative development; used via GitHub web interface and VSCode.
\end{itemize}

The author identified scope, validated all content against official Falcon BMS documentation, conducted final text review, made all organizational, editorial, and publication decisions, and bears full responsibility for the guide's accuracy and presentation.

%--------------------------------------------------------------------------
% SUBSECTION 1.4.3: LICENSE AND CANONICAL SOURCE
%--------------------------------------------------------------------------
\subsection{License and canonical source}
\label{sec:C1-S4-S3}

The canonical source for this guide is the project repository:

\begin{center}
	\url{https://github.com/carlos-nader/tms-dms-cms-usage-guide}
\end{center}

The repository contains the authoritative version of all project files, including governance documents and version history. Translations, adaptations, and derivative works should reference this repository as the original source.

The guide is also available as a web-readable version at:

\begin{center}
	\url{https://tms-dms-cms-usage-guide.readthedocs.io}
\end{center}

The web version provides navigable, searchable access to all guide content and is recommended for online reference alongside the simulator.

This document is released under the \textbf{Creative Commons Attribution-NonCommercial 4.0 International} (CC BY-NC 4.0) license. Sharing, copying, translating, and adapting are permitted with attribution. \textbf{Commercial use is prohibited}. Derivative works must maintain the same license and cannot claim official status. Full legal terms: \url{https://creativecommons.org/licenses/by-nc/4.0/}

\paragraph{Attribution:}
When sharing, translating, or adapting, please include: title, author (Carlos ``Metal'' Nader), source repository URL, and license.

%--------------------------------------------------------------------------
% SUBSECTION 1.4.4: DISCLAIMER
%--------------------------------------------------------------------------
\subsection{Disclaimer}
\label{sec:C1-S4-S4}

This is an unofficial, community-made document, provided ``as-is'' for \textbf{educational and simulation training} purposes only.

No copyrighted material is reproduced directly.

%============================================================================
% CHAPTER 2 — HOTAS FUNDAMENTALS
%============================================================================
\chapter{HOTAS Fundamentals}
\label{chap:C2}
The F-16 HOTAS (Hands On Throttle And Stick) controls consist of switches on the throttle and side-stick controller. The Display Management Switch (DMS), Target Management Switch (TMS), and Countermeasures Management Switch (CMS) --- the three HOTAS switches covered in this guide --- control display management, sensor management, and defensive systems.

SOI and Master Mode determine how HOTAS inputs are interpreted:

\begin{itemize}
    \item \textbf{Sensor of Interest (SOI):} determines which display receives HOTAS inputs.
    \item \textbf{Master Modes (NAV, A-A, A-G):} define the available weapon systems and sensor configurations.
\end{itemize}

For the complete HOTAS control reference, see \dashrefs{2.1.1.1.4 and 2.1.5}.

%============================================================================
% SECTION 2.1 — SENSOR OF INTEREST (SOI)
%============================================================================
\section{Sensor of Interest (SOI)}
\label{sec:C2-S1}

The F-16 cockpit contains three displays: the Head-Up Display (HUD) and two Multifunction Displays (MFD). The Sensor of Interest (SOI) determines which of these displays currently receives HOTAS inputs. At any moment, \textit{only one} display can be designated as SOI. The active format on the display determines how the input is interpreted.

%--------------------------------------------------------------------------
% SUBSECTION 2.1.1 — SOI FUNDAMENTALS
%--------------------------------------------------------------------------
\subsection{SOI Fundamentals}
\label{sec:C2-S1-S1}

The SOI mechanism manages \textit{where} HOTAS inputs are directed, not \textit{what} those inputs do. The specific function of a HOTAS input --- for example, whether \textbf{TMS} Up designates a target, breaks a track, or cycles through a mode --- depends on delivery modes, current Master Mode, and active sensor format.

%--------------------------------------------------------------------------
% SUBSECTION 2.1.2 — SOI SYMBOLOGY
%--------------------------------------------------------------------------
\subsection{SOI Symbology}
\label{sec:C2-S1-S2}

SOI designation is indicated by a visual marker on each display:

\begin{itemize}
	\item \textbf{HUD as SOI:} An asterisk symbol (\texttt{*}) appears in the upper left corner of the HUD, above the airspeed scale.
	\item \textbf{MFD as SOI:} A border outline is drawn around the edges of the MFD display, forming a box that distinguishes the SOI display from the non-SOI display.
	\item \textbf{NOT SOI indication:} The text \texttt{NOT SOI} appears in the center of any SOI-capable MFD format that is not currently designated as SOI. This indicator does not appear on MFD formats that cannot be designated as SOI.
\end{itemize}

%--------------------------------------------------------------------------
% SUBSECTION 2.1.3 — SOI-CAPABLE DISPLAYS
%--------------------------------------------------------------------------
\subsection{SOI-Capable Displays}
\label{sec:C2-S1-S3}

SOI designation depends on Master Mode and MFD format.

Master Mode determines SOI eligibility. MFD formats can be designated as SOI in \textit{all Master Modes}. The HUD, however, can be designated as SOI \textit{only in NAV and A-G Master Modes}.

MFD format determines SOI eligibility \textit{within} the MFD. Sensor formats that accept targeting and cursor control inputs --- such as FCR, TGP, and WPN --- can be designated as SOI. Informational and configuration formats, however --- such as DTE, TCN, and FLCS --- cannot be designated as SOI.

For SOI availability by Master Mode and MFD format, see \secref{sec:C3-S1}.

%============================================================================
% SECTION 2.2 — MASTER MODES
%============================================================================
\section{Master Modes}
\label{sec:C2-S2}

Master Mode defines the active sensor, display, and weapon configuration of the F-16 avionics suite. It also determines SOI availability and switch functions.

%--------------------------------------------------------------------------
% SUBSECTION 2.2.1 — PRIMARY AND OVERRIDE MASTER MODES
%--------------------------------------------------------------------------
\subsection{Primary and Override Master Modes}
\label{sec:C2-S2-S1}

Three primary Master Modes (NAV, A-A, and A-G) and two override modes (DGFT and MSL OVRD) are available in the F-16:

\begin{itemize}

	\item \textbf{Navigation (NAV):} The default Master Mode, active when no other mode is selected. Both air-to-air and air-to-ground sensor modes are available in NAV. Navigation symbology is displayed on the HUD.
    \item \textbf{Air-to-Air (A-A):} Selected via the A-A button on the ICP. Only air-to-air sensor modes are available in A-A. Air-to-air weapon selection and engagement symbology are displayed on the HUD.
    \item \textbf{Air-to-Ground (A-G):} Selected via the A-G button on the ICP. Both air-to-air and air-to-ground sensor modes are available in A-G. Air-to-ground weapon delivery symbology is displayed on the HUD.
    \item \textbf{Dogfight (DGFT) and Missile Override (MSL OVRD):} Selected via the Dogfight/Missile Override switch on the throttle — DGFT outboard, MSL OVRD inboard. Only air-to-air sensor modes are available in both modes. DGFT and MSL OVRD take precedence over all other Master Modes except Emergency Jettison.

\end{itemize}

HOTAS switch inputs are interpreted within the active master mode configuration --- sensor, display, and weapon options ---, which can be pre-configured via the Data Transfer Cartridge (DTC) or adjusted manually in-flight. For complete Master Mode details, see \dashref{2.1.1.2.1}.

%============================================================================
% SECTION 2.3 — DMS, TMS, AND CMS
%============================================================================
\section{DMS, TMS, and CMS}
\label{sec:C2-S3}

DMS, TMS, and CMS — including all HOTAS tables — are covered in \chapref{chap:C3}, \chapref{chap:C4}, and \chapref{chap:C5}, respectively.

%--------------------------------------------------------------------------
% SUBSECTION 2.3.1 — DISPLAY MANAGEMENT SWITCH (DMS)
%--------------------------------------------------------------------------
\subsection{Display Management Switch (DMS)}
\label{sec:C2-S3-S1}

The Display Management Switch (DMS) is a four-way hat switch on the side-stick controller. Its primary functions are SOI designation and MFD format cycling (see \dashref{2.1.5}). The DMS selects \textit{which display or sensor} receives HOTAS inputs. It does not designate targets or control sensor modes.

%--------------------------------------------------------------------------
% SUBSECTION 2.3.2 — TARGET MANAGEMENT SWITCH (TMS)
%--------------------------------------------------------------------------
\subsection{Target Management Switch (TMS)}
\label{sec:C2-S3-S2}

The Target Management Switch (TMS) is a four-way hat switch on the side-stick controller. Its primary functions are target designation, sensor management, and markpoint creation (see \dashref{2.1.5}). The function of each TMS direction depends on the Master Mode, the active sensor mode, and the current SOI display.

%--------------------------------------------------------------------------
% SUBSECTION 2.3.3 — COUNTERMEASURES MANAGEMENT SWITCH (CMS)
%--------------------------------------------------------------------------
\subsection{Countermeasures Management Switch (CMS)}
\label{sec:C2-S3-S3}

The Countermeasures Management Switch (CMS) is a four-way hat switch on the side-stick controller. Its primary functions are countermeasures dispensing and ECM control (see \dashref{2.1.5}). CMS actuation is independent of both SOI designation and Master Mode.

%============================================================================
% CHAPTER 5: CMS — Countermeasures Management Switch
%============================================================================
\chapter{CMS – Countermeasures Management Switch}
\label{chap:C5}

%============================================================================
% SECTION 5.1: Concept
%============================================================================
\section{Concept}
\label{sec:C5-S1}

The Countermeasures Management Switch (CMS) is a four-way hat switch mounted on the side-stick controller. It controls the aircraft's defensive avionics:

\begin{itemize}
  \item ALE-47 CMDS (Countermeasures Dispenser Set) --- the chaff and flare dispensing system; and
  \item ECM (Electronic Countermeasures) --- the radar-jamming system.
\end{itemize}

ECM configuration depends on block and variant: either external pods (ALQ-131/ALQ-184) or internal avionics (IDIAS).

CMS actuation procedures are presented in \secref{sec:C5-S2} and block and variant differences in \secref{sec:C5-S3}.

For complete CMDS (including chaff and flare dispensing programs) and ECM operation, see \dashrefs{2.7.1 and 2.7.2}, respectively. Diagrams of the side-stick controller and throttle functions are provided in \dashref{2.1.5}. \figref{fig:f16_hotas_cms_location} below shows the CMS location on the F-16 side-stick controller.

\begin{figure}[H]
\centering
\includegraphics[width=0.65\textwidth]{F-16_Side_Stick_Controller-1.jpg}
\caption{F-16 Throttle and Side-Stick Controller. Image by Falconpedia (\url{falcon4.wikidot.com}), via Wikimedia Commons (\url{https://commons.wikimedia.org/wiki/File:F-16_Side_Stick_Controller.jpg}), licensed under the Creative Commons Attribution-Share Alike 3.0 Unported (CC BY-SA 3.0) license.}
\label{fig:f16_hotas_cms_location}
\end{figure}

%============================================================================
% SECTION 5.2: CMS Actuation
%============================================================================
\section{CMS Actuation}
\label{sec:C5-S2}

The tables below cover CMS actuation organized by system: CMDS (Manual, Automatic, Semi-Automatic, Bypass, and Standby modes) and ECM (external pods ALQ-131/ALQ-184 and internal IDIAS).

All CMS actuation described in this chapter is \textit{independent} of the Master Mode (NAV, A-A or A-G) selected.

Table columns are defined as follows:

\begin{itemize}
  \item \textbf{Mode}: The operational context (Master Mode).
  \item \textbf{Direction}: CMS hat direction (Up, Down, Left, Right).
  \item \textbf{Action}: Press type (Short, Long, Long Hold).
  \item \textbf{Function}: What the CMS command activates or controls.
  \item \textbf{Effect / Nuance}: Resulting system behavior, including tactics and constraints.
  \item \textbf{Dash34}: Reference section in the Dash-34 manual.
  \item \textbf{Training}: Recommended BMS training mission(s) (see \trnman{} for descriptions).
\end{itemize}

%============================================================================
% SECTION 5.2.1: CMDS
%============================================================================
\subsection{CMDS}
\label{sec:C5-S2-S1}

The ALE-47 CMDS operates in five modes selected via the CMDS CCU MODE knob: Manual (MAN), Automatic (AUTO), Semi-Automatic (SEMI), Bypass (BYP), and Standby (STBY).

The chaff and flare dispensing program executed in MAN, AUTO, and SEMI modes --- Programs 1--4 --- is selected via the PRGM knob on the CMDS panel. The PRGM knob has no effect in BYP or STBY.

Only Programs 1--4 are selectable. Programs 5 and 6 are independent of the PRGM knob position:

\begin{itemize}
  \item Program 5 is activated via the cockpit slap switch --- outside the scope of CMS --- and is not covered in this guide; see \dashref{2.7.2.2} for details.
  \item Program 6 is a special-purpose program always available for manual dispensing, independent of the PRGM knob position.
\end{itemize}

%============================================================================
% SECTION 5.2.1.1: Manual Mode
%============================================================================
\subsubsection{Manual Mode}
\label{sec:C5-S2-S1-S1}

In Manual (MAN) mode, the pilot executes countermeasure programs (1--4 and 6) via \textbf{CMS Up} and \textbf{CMS Left}.

\begin{hotastable}{CMS Actuation with CMDS Manual Mode}\label{tab:C5-S2-S1-S1-CMDS-Manual}
CMDS MAN & Up & Short & Execute Program 1--4 &
Executes the program selected via the CMDS panel PRGM knob once per press. No threat sensing; no automatic triggering. Overrides any AUTO dispensing, if running.
& 2.7.2.2 & 18, 28 \\
CMDS MAN & Left & Short & Execute Program 6 &
Program 6 is dispensed independently of the PRGM knob position.
& 2.7.2.2 & 19 \\
\end{hotastable}

%============================================================================
% SECTION 5.2.1.2: Automatic Mode
%============================================================================
\subsubsection{Automatic Mode}
\label{sec:C5-S2-S1-S2}

In Automatic (AUTO) mode, the ALE-47 CMDS dispenses the selected program (1–4) \textit{continuously} in response to threats detected by the Radar Warning Receiver (RWR).

A \textbf{CMS Down} press grants consent for AUTO dispensing to start. It will continue until the pilot presses \textbf{CMS Right} or expendables are exhausted.

Even in AUTO mode, manual dispensing is still available.

\begin{hotastable}{CMS Actuation with CMDS Automatic Mode}\label{tab:C5-S2-S1-S2-CMDS-Auto}
CMDS AUTO & Up & Short & Execute Program 1--4 &
Executes the program selected via the CMDS panel PRGM knob once, interrupting any ongoing AUTO dispensing. After execution, AUTO resumes if the threat persists.
& 2.7.2.2 & 18, 28 \\
CMDS AUTO & Left & Short & Execute Program 6 &
Program 6 is dispensed independently of the PRGM knob position. Overrides AUTO; after execution, AUTO resumes.
& 2.7.2.2 & 19 \\
CMDS AUTO & Down & Short & Give Consent; Enable AUTO Dispensing &
Grants consent for AUTO dispensing. RWR-detected threats trigger automatic execution of the program (1--4) selected via the CMDS panel PRGM knob. Dispensing continues until the threat clears, the pilot presses CMS Right, or expendables are exhausted. Consent state persists if the pilot switches to MAN mode. If the pilot re-engages AUTO, dispensing resumes immediately on the next detected threat.
& 2.7.2.1 & 18, 28 \\
CMDS AUTO & Right & Short & Cancel Consent; Disable AUTO &
Removes CMDS consent and places the ALE-47 in Standby. Automatic dispensing halts immediately. The pilot must press CMS Down again to resume AUTO operation.
& 2.7.2.1 & 18, 28 \\
\end{hotastable}

%============================================================================
% SECTION 5.2.1.3: Semi-Automatic Mode
%============================================================================
\subsubsection{Semi-Automatic Mode}
\label{sec:C5-S2-S1-S3}

In Semi-Automatic (SEMI) mode, the CMDS control unit displays \texttt{DISPENSE RDY} and sounds a \texttt{COUNTER} voice message every time a threat is detected. For each \texttt{COUNTER} voice message, a \textbf{CMS Down} press grants consent and the ALE-47 CMDS then dispenses \textit{one selected program}.

Even in SEMI mode, manual dispensing is still available.

\begin{hotastable}{CMS Actuation with CMDS Semi-Automatic Mode}\label{tab:C5-S2-S1-S3-CMDS-Semi}
CMDS SEMI & Up & Short & Execute Program 1--4 &
Executes the program selected via the CMDS panel PRGM knob once, independent of RWR threat state. After execution, the CMDS returns to standby and resumes threat detection.
& 2.7.2.2 & 18, 28 \\
CMDS SEMI & Left & Short & Execute Program 6 &
Program 6 is dispensed independently of the PRGM knob position. After execution, the CMDS returns to standby and resumes threat detection.
& 2.7.2.2 & 19 \\
CMDS SEMI & Down & Short & Give Consent; Dispense One Program &
Dispenses one program (1--4) in response to a \texttt{COUNTER} voice message. Consent state is retained: if the pilot switches to AUTO, the CMDS begins dispensing immediately on the next detected threat.
& 2.7.2.1 & 18, 28 \\
CMDS SEMI & Right & Short & Cancel Consent; Return to Standby &
Removes SEMI consent and places the CMDS in Standby. \texttt{COUNTER} voice messages cease.
& 2.7.2.1 & 18, 28 \\
\end{hotastable}

%============================================================================
% SECTION 5.2.1.4: Bypass and Standby Modes
%============================================================================
\subsubsection{Bypass and Standby Modes}
\label{sec:C5-S2-S1-S4}

In BYP and STBY, both \textit{non-combat modes}, AUTO and SEMI dispensing are not available. Consent state is not tracked in either mode.

\paragraph{Bypass (BYP):}
In BYP, the CMDS program selection is bypassed: \textbf{CMS Up} dispenses \textit{exactly} one chaff and one flare per press. Program 6 remains available via \textbf{CMS Left}.

\begin{hotastable}{CMS Actuation with CMDS Bypass Mode}\label{tab:C5-S2-S1-S4-CMDS-BYP}
CMDS BYP & Up & Short & Dispense One Chaff + One Flare &
The PRGM knob setting has no effect. AUTO and SEMI dispensing are not available in BYP.
& 2.7.2.1 & --- \\
CMDS BYP & Left & Short & Execute Program 6 &
Program 6 is dispensed independently of the PRGM knob position.
& 2.7.2.2 & --- \\
\end{hotastable}

\paragraph{Standby (STBY):}
In STBY, CMDS programs and bingo quantities can be edited via the UFC. This mode inhibits all chaff and flare release: \textit{no} CMS direction produces a dispensing effect.
For program editing procedures, see \dashref{2.7.2.1 and 2.7.2.3}.

%============================================================================
% SECTION 5.2.2: ECM
%============================================================================
\subsection{ECM}
\label{sec:C5-S2-S2}

The ECM system operates in two configurations:

\begin{itemize}
  \item External pods (ALQ-131, ALQ-184) --- controlled via \textbf{CMS Down}; and
  \item Internal IDIAS --- controlled via \textbf{CMS Left}.
\end{itemize}

%============================================================================
% SECTION 5.2.2.1: External ECM Pod (ALQ-131/ALQ-184)
%============================================================================
\subsubsection{External ECM Pod (ALQ-131/ALQ-184)}
\label{sec:C5-S2-S2-S1}

For the external ECM pod to operate, its power switch must be in OPR. This enables the XMIT switch on the ECM control panel to set the operation mode:

\begin{itemize}
  \item XMIT 1: Avionics Priority, AFT antenna only
  \item XMIT 2: ECM Priority, both FWD and AFT antennas
  \item XMIT 3: Active Jam, continuous transmission
\end{itemize}

After setting the operation mode, \textbf{CMS Down} grants consent for ECM pod transmission. The ECM Enable light on the left miscellaneous panel illuminates.

\begin{hotastable}{CMS Actuation with External ECM Pod}\label{tab:C5-S2-S2-S1-ECM-Pod}
ECM Pod & Down & Short & Grant ECM Consent &
The pod transmits in the operation mode set by the XMIT switch. Transmission continues until the pilot presses CMS Right or the RF switch is moved out of NORM.
& 2.7.4.1.1, 2.7.4.2.5 & 28 \\
ECM Pod & Right & Short & Remove ECM Consent &
ECM consent is removed. The pod is placed in Standby and the ECM Enable light extinguishes. The pilot must press CMS Down again to resume transmission.
& 2.7.4.1.1, 2.7.4.2.5 & 28 \\
\end{hotastable}

%============================================================================
% SECTION 5.2.2.2: Internal ECM (IDIAS)
%============================================================================
\subsubsection{Internal ECM (IDIAS)}
\label{sec:C5-S2-S2-S2}

For IDIAS to operate, its power switch must be in ON. Then, with the XMTR switch in OPER, \textbf{CMS Left} cycles the ECM operation mode.

The active operation mode is displayed on the HUD, FCR, RWR, and TFR pages.

\begin{hotastable}{CMS Actuation with Internal ECM (IDIAS)}\label{tab:C5-S2-S2-S2-IDIAS}
IDIAS ECM & Left & Short & Cycle ECM Operation Mode &
The IDIAS mode advances with each short press: \texttt{STBY} $\rightarrow$ \texttt{AVNC} (Avionics Priority) $\rightarrow$ \texttt{ECM} (ECM Priority) $\rightarrow$ \texttt{AVNC} $\rightarrow$ \texttt{ECM}.
& 2.7.4.1.2 & 28 \\
IDIAS ECM & Right & Short & Set ECM to Standby &
IDIAS is set to \texttt{STBY}. The pilot must press \textbf{CMS Left} again to return to \texttt{AVNC} or \texttt{ECM}.
& 2.7.4.1.2 & 28 \\
\end{hotastable}

%============================================================================
% SECTION 5.2.3: CMDS and ECM Joint Consent
%============================================================================
\subsection{CMDS and ECM Joint Consent}
\label{sec:C5-S2-S3}

On aircraft equipped with an \textit{external} ECM pod, \textbf{CMS Down} grants consent to \textit{both} the ALE-47 CMDS and the external ECM pod. This joint consent applies to CMDS AUTO and SEMI modes \textit{only}.

As \textbf{CMS Down} has no function on IDIAS-equipped aircraft, the joint consent does not apply.

\begin{hotastable}{CMS Interaction with CMDS and ECM (Consent and Constraints)}\label{tab:C5-S2-S3-CMDS-ECM}
CMDS AUTO SEMI + ECM Pod & Down & Short & Joint Consent (CMDS + ECM) &
Simultaneous consent to CMDS and the external ECM pod.
& 2.7.2.1, 2.7.4.2.5 & 18, 28 \\
\end{hotastable}

%============================================================================
% SECTION 5.2.4: Constraints and Edge Cases
%============================================================================
\subsection{Constraints and Edge Cases}
\label{sec:C5-S2-S4}

\paragraph{State Tracking:}

In CMDS AUTO and SEMI modes, consent state is \textit{not cleared} when MAN mode is selected. If the pilot re-engages AUTO or SEMI without pressing \textbf{CMS Down}, the CMDS begins dispensing immediately on the next detected threat. As BYP and STBY modes do not have a consent state, switching to either mode does not retain any prior consent.

\paragraph{Bingo Quantity Behavior:}

In AUTO mode, the CMDS continues dispensing when expendables fall to or below the bingo quantity. The pilot must press \textbf{CMS Right} to disable AUTO dispensing, or switch to MAN mode.

\paragraph{RF Switch Override:}

When ECM consent is active, the ECM system is placed in Standby if the RF switch is moved from NORM to QUIET or SILENT. Consent is \textit{not} restored when the RF switch is returned to NORM. The pilot must press \textbf{CMS Down} again to resume ECM transmission.

\paragraph{Ground Safety:}

On the ground, if the pilot holds \textbf{CMS Down}, the ECM pod may radiate. Ground personnel must be clear of the radiation area.

%============================================================================
% SECTION 5.3: Block and Variant Configurations
%============================================================================
\section{Block and Variant Configurations}
\label{sec:C5-S3}

CMS interaction with CMDS is \textit{uniform} across all F-16 block and variant configurations (see \secref{sec:C5-S2-S1}).

F-16 block and variant configurations with external pod ECM, however, use \textit{different} CMS procedures and ECM control panels (see \secref{sec:C5-S2-S2}):

\begin{table}[H]
\centering
\caption{ECM Configuration by Block and Variant}
\label{tab:C5-S3-ecm-config-summary}
\small
\renewcommand{\arraystretch}{1.25}
\begin{tabular}{|>{\raggedright\arraybackslash}p{3.6cm}|>{\raggedright\arraybackslash}p{3.6cm}|>{\raggedright\arraybackslash}p{4.6cm}|}
\hline
\rowcolor{headerblue}
\textcolor{white}{\textbf{ECM Configuration}} & \textcolor{white}{\textbf{CMS Direction}} & \textcolor{white}{\textbf{Panel Control}} \\
\hline
External Pod (ALQ-131/ALQ-184) & CMS Down (consent) & XMIT switch (positions: 1, 2, 3) \\
\hline
Internal IDIAS & CMS Left (mode cycling) & XMTR switch (positions: STBY, OPER) \\
\hline
\end{tabular}
\end{table}

%============================================================================
% SECTION 5.3.1: ECM Configurations by Block and Variant
%============================================================================
\subsection{ECM Configurations by Block and Variant}
\label{sec:C5-S3-S1}

The block and variant configurations listed in the following tables are sourced from Dash-34 documentation. Not all may be available as selectable aircraft in Falcon BMS 4.38.1.

%============================================================================
% SECTION 5.3.1.1: ECM Pods (ALQ-131/ALQ-184)
%============================================================================
\subsubsection{ECM Pods (ALQ-131/ALQ-184)}
\label{sec:C5-S3-S1-S1}

These block and variant configurations use the ECM Pod Control Panel with manual jamming program selection. For detailed procedures, see \secref{sec:C5-S2-S2-S1}.

\begin{table}[H]
\centering
\caption{External ECM Pod Blocks and Variants}
\label{tab:C5-S3-external-ecm-pods}
\small
\renewcommand{\arraystretch}{1.25}
\begin{tabular}{|>{\raggedright\arraybackslash}p{3.0cm}|>{\raggedright\arraybackslash}p{5.8cm}|}
\hline
\rowcolor{headerblue}
\textcolor{white}{\textbf{Operator}} & \textcolor{white}{\textbf{Block/Variant}} \\
\hline
USAF & Blocks 40/42/50/52 (CM designation) \\
\hline
NATO & Block 15 operators (Belgium, Denmark, Netherlands, Norway) \\
\hline
International & Egypt, Korea KF-16C Block 32 \\
\hline
\end{tabular}
\end{table}

%============================================================================
% SECTION 5.3.1.2: Internal ECM (IDIAS)
%============================================================================
\subsubsection{Internal ECM (IDIAS)}
\label{sec:C5-S3-S1-S2}

These block and variant configurations use the IDIAS ECM control panel. For detailed procedures, see \secref{sec:C5-S2-S2-S2}.

\begin{table}[H]
\centering
\caption{Internal ECM (IDIAS) Blocks and Variants}
\label{tab:C5-S3-internal-ecm}
\small
\renewcommand{\arraystretch}{1.25}
\begin{tabular}{|>{\raggedright\arraybackslash}p{3.0cm}|>{\raggedright\arraybackslash}p{5.8cm}|}
\hline
\rowcolor{headerblue}
\textcolor{white}{\textbf{Operator}} & \textcolor{white}{\textbf{Block/Variant}} \\
\hline
Israel & F-16I Sufa Block 52, Barak I Block 30, Barak II Block 40 \\
\hline
International & Greek HAF (Blocks 50 PXII, 52 PXIII, 52+ PXIV Advanced),
Korea KF-16C Block 52, Singapore F-16D Block 52 \\
\hline
\end{tabular}
\end{table}

%============================================================================
% END OF SECTION
%============================================================================

\end{document}